\documentclass[10pt,a4paper]{amsart}
\usepackage[margin=19mm]{geometry}
\usepackage{amsmath,amssymb,mathtools}
\usepackage[scr=rsfs,cal=euler]{mathalfa}
\usepackage{xfrac}
\usepackage[unicode,hidelinks,bookmarks=false]{hyperref}
\newcommand{\ie}{i.e.\ }
\newcommand{\cf}{cf.\ }
\newcommand{\resp}{respectively\ }
\newcommand{\Subsection}{\subsection}
\providecommand{\nobreakdash}{\nobreak\hbox{-}}
\setlength{\emergencystretch}{2em}
\multlinegap0pt
\usepackage{bm}
\usepackage{bbm}
\usepackage{dsfont}
\usepackage{xspace}
\usepackage{cancel}
\usepackage{mathtools}
\usepackage{amsthm}
\makeatletter
\@ifundefined{lem}{\newtheorem{lem}{Lemma}}{}
\makeatother
\makeatletter
\expandafter\ifx\csname solucion\endcsname\relax
\newenvironment{solucion}{\noindent \emph{Solución.}}{}
\fi
\makeatother
\title[On the Borel complexity of the space of left-orderings of ni]{On the Borel complexity of the space of left-orderings of nilpotent groups}
\author{Emir Molina Tauc\'an}
\date{Source: arXiv:2507.06953v1 (CC BY 4.0)}
\begin{document}
\maketitle

\noindent\emph{Source and licence.} On the Borel complexity of the space of left-orderings of nilpotent groups. Version arXiv:2507.06953v1 (CC BY 4.0), distributed under the Creative Commons Attribution 4.0 International licence, as recorded in the arXiv licence metadata for this exact version and shown on the abstract page https://arxiv.org/abs/2507.06953v1. Full rights notice: licenses/arxiv-2507.06953-CC-BY-4.0.txt.\par\smallskip

\noindent\textit{Editorial scope.} Counted unit: the Lem whose source label is \texttt{lemacono} in the article (source0.tex lines 394--398, source label \texttt{lemacono}), delivered together with the article's own complete original proof of it (source0.tex lines 399--438, one \texttt{proof} environment of 40 lines). No mathematics is written, repaired or re-derived: statement and proof are the article's own text line for line. Required context retained uncounted: none. Every definition, display and label of those lines is retained unchanged. Omitted from this package and not redistributed: 1--393, 439--883. Nothing inside a retained excerpt is omitted or altered: every retained body is delivered exactly as the article's lines are written, so no omission receipt of any kind accompanies this package. Citations: the retained excerpts carry no citation command at all, so no bibliography entry of the article is transcribed and no reference is redistributed. Reference adaptation: none was needed and none was made; every reference inside the delivered unit resolves inside it, so no unresolved reference or citation is produced. Printed slips kept literal: no printed word, symbol, hypothesis or number of the counted statement or of its proof was altered. Layout and class: the article's own class is \texttt{article} with packages including babel, graphicx, amsmath, mathtools, amsthm, amssymb, float, geometry, enumerate, ragged2e, array, tikz, vmargin, marvosym; the wrapper is the lane's pinned \texttt{amsart} editorial wrapper at 10pt on A4, which restates the theorem-like environments the retained text needs and adds the article's own macro definitions that the retained text needs (none), with extra wrapper packages (bm, bbm, dsfont, xspace, cancel, mathtools). The delivered numbering is the unit's own. Assets: no retained span contains a figure environment, a graphics-inclusion command, a PDF-page inclusion, a table or a TikZ picture; no image file is copied into this package, the package declares no asset dependency and no image file is redistributed with it. Licence: version arXiv:2507.06953v1 of this article is distributed under the Creative Commons Attribution 4.0 International licence, as recorded in the arXiv licence metadata for this exact version and shown on the abstract page https://arxiv.org/abs/2507.06953v1. A full rights notice is delivered at licenses/arxiv-2507.06953-CC-BY-4.0.txt. Characters: every retained excerpt is pure ASCII, so the delivered engine, XeLaTeX with its default Latin Modern text font, reports no missing character.
\begin{lem}
\label{lemacono}
    Let $\vec{v}=(x_1,\dots,x_n)$ be a vector in $\mathbb{R}^n$, where $x_n>0$, and $A$ a matrix in $L_n\setminus\{I_n\}$.
    Then there exists an integer vector $\vec{w}\in \mathbb{Z}^n$ such that $\langle \vec{v} , \vec{w} \rangle>0$ but $\langle (A^{-1})^t(\vec{v}) , \vec{w} \rangle<0.$
\end{lem}

\begin{proof}
Let us denote by $H_{v}$ the orthogonal complement of $\langle\ \vec{v}\ \rangle$ and define $\vec{w_i}:= \text{sgn}(a_i)\cdot(\vec{e_i}-\frac{x_i}{x_n}\vec{e}_n)$ for all $i=1,\dots,n-1$. 
Here, we use the convention sgn$(0)=1$. 

The reader can verify that $\{\vec{w}_1,\dots, \vec{w}_{n-1}\}$ is a $\mathbb{R}$-linearly independent set and, therefore, $H_v=\langle\ \vec{w}_1,\dots, \vec{w}_{n-1}\ \rangle$.
Moreover, given that $\langle\vec{v},\vec{x}\rangle=\langle A^t(A^{-1})^t(\vec{v}),\vec{x}\rangle =\langle (A^{-1})^t(\vec{v}),A(\vec{x})\rangle$ for all $\vec{x}\in \mathbb{R}^n$, we also conclude that $H_{(A^{-1})^tv}=\langle A(\vec{w}_1),\dots, A(\vec{w}_{n-1})\rangle.$

A simple calculus shows that $A(\vec{w}_i)=\vec{w}_i+|a_i|\cdot \vec{e_n}$. Thus $\langle \vec{v},A(\vec{w}_i)\rangle=\langle\vec{v},\vec{w}_i\rangle +|a_i|\langle \vec{v},\vec{e_n}\rangle=|a_i|x_n\geq 0$ for all indices $i$.
Of course, we have that $\langle\vec{v},A(\vec{w}_i)\rangle>0$ when $a_i\neq 0.$

Consider now $i_0$ the smallest index $i$ such that $a_i\neq 0$ and fix the vector $A(\vec{w}_{i_0})=\vec{w}_{i_0}+|a_{i_0}|\cdot\vec{e}_n$. We also consider the convex cone $C$ generated by $\vec{w}_1,\dots, \vec{w}_{i_0},A(\vec{w}_{i_0}),\dots, \vec{w}_{n-1}$. Recall that $C=\{\ \sum_{i=1}^{n-1}\alpha_i\cdot\vec{w}_i  + \beta_{i_0}A(\vec{w}_{i_0})\ |\ \alpha_{i},\beta_{i_0}\in \mathbb{R}_{\geq 0}\ \}$.

 To conclude with our proof, we need the following:\\ 

    \noindent\textit{Claim.} If $B=\{\vec{v}_1,\dots, \vec{v}_n\}$ is a basis of the vector space $\mathbb{R}^n$ and $C$ the convex cone generated by the same elements, then Int$(C)\cap \mathbb{Z}^n\neq \emptyset$.\\

    \noindent\textit{Proof of Claim.} Let us consider the convex cone $C_0$ generated by the canonical basis $B_0$. We also consider the change-of-basis matrix $A$ of $B_0$ to $B$. The reader can note that $A: \mathbb{R}^n\rightarrow \mathbb{R}^n$ is a homeomorphism such that $A(C_0)=C$. 
    
    Due to Int$(C_0)$ being trivially nonempty and $A$ being a homeomorphism, we conclude that Int$(C)$ is also nonempty. Then there exist $\vec{w}\in C$ and $\epsilon>0$ such that $B(\vec{w},\epsilon)\subseteq C$, where $B(\vec{w},\epsilon)$ denoted the open ball with radius $\epsilon$ (with respect to the Euclidean metric in $\mathbb{R}^n$) and center $\vec{w}$. Moreover, $B(\lambda\vec{w},\lambda\epsilon)=\lambda B(\vec{w},\epsilon)\subseteq C$ for all $\lambda>0$.
    
    Finally, there must exist $\lambda>0$ large enough such that $B(\lambda\vec{w},\lambda\epsilon)\cap \mathbb{Z}^n\neq \emptyset$. $\hfill\square$\\

A quick thought will convince us that $\{\vec{w}_1,\dots, \vec{w}_{i_0},A(\vec{w}_{i_0}),\dots, \vec{w}_{n-1}\}$ is a basis of $\mathbb{R}^n$. From the previous claim, it follows that there exists a integer vector $\vec{z}\in \text{Int}(C)$.

In general, if $\vec{x}=\sum_{i=1}^{n-1}\alpha_i\cdot\vec{w}_i  + \beta_{i_0}A(\vec{w}_{i_0})\in C$, then
\begin{equation}   
\label{igual1}
\langle\vec{v},\vec{x}\rangle=\sum_{i=1}^{n-1}\alpha_{i}\cdot\langle\vec{v},\vec{w}_i \rangle + \beta_{i_0}\cdot \langle \vec{v},A(\vec{w}_{i_0})\rangle=\beta_{i_0}\cdot \langle \vec{v},A(\vec{w}_{i_0})\rangle=\beta_{i_0}x_n\geq 0,
\end{equation}
\begin{equation}
\label{igual2}
\text{and }\ \langle (A^{-1})^t(\vec{v}),\vec{x}\rangle=\sum_{i=1}^{n-1}\alpha_i\cdot\langle (A^{-1})^t(\vec{v}),\vec{w}_i\rangle=\sum_{i=1}^{n-1}\alpha_i\cdot\langle \vec{v},(A^{-1})(\vec{w}_i)\rangle=-\sum_{i=1}^{n-1}\alpha_i|a_{i}|x_n\leq 0.    
\end{equation}

Given that $\vec{z}\in \text{Int}(C)$, then $\beta_{i_0}$ there must be a positive number. In fact, if $\beta_{i_0}=0$, then $\vec{x}\in H_v$ and, therefore,  $\langle\vec{v},\vec{x}\rangle=0$. On the other hand, there exists $\epsilon>0$ such that $\vec{x}-\epsilon\cdot\vec{v}\in C$. Finally, $\langle\vec{v},\vec{x}-\epsilon\cdot\vec{v}\rangle=-\epsilon\cdot\langle\vec{v},\vec{v}\rangle<0$. This inequality contradicts equation (\ref{igual1}).

Similarly, there must exist some index $i$ such that $\alpha_i> 0$ and $a_i\neq 0$. In fact, if $\alpha_i=0$ for all indices $i$ such that $a_i\neq 0$, then $\vec{x}\in A(H_v)=H_{(A^{-1})^tv}$. Moreover, there exists $\epsilon>0$ such that $\vec{x}+\epsilon\cdot(A^{-1})^t(\vec{v})\in C$. Finally, $\langle (A^{-1})^t(\vec{v}), \vec{x}+ \epsilon\cdot(A^{-1})^t(\vec{v})\rangle=\epsilon\cdot\langle (A^{-1})^t(\vec{v}), (A^{-1})^t(\vec{v})\rangle>0$. This inequality contradicts equation (\ref{igual2}).

Finally, we conclude that $\vec{z}$ satisfies $\langle\vec{v},\vec{z}\rangle>0$ and $\langle(A^{-1})^t(\vec{v}),\vec{z}\rangle<0.$ 
\end{proof}

\end{document}
